\documentclass[UTF8,a4paper,11pt]{ctexart}
\usepackage[a4paper,margin=2.3cm]{geometry}
\usepackage{amsmath,amssymb,mathtools}
\usepackage{enumitem}
\usepackage[hidelinks]{hyperref}
\usepackage{fancyhdr}
\setlist[itemize]{leftmargin=2.2em,itemsep=0.25em,topsep=0.35em}
\setlength{\parindent}{2em}\setlength{\parskip}{0.35em}
\pagestyle{fancy}\fancyhf{}\fancyhead[L]{离散数学与结构（I）}\fancyhead[R]{Lec 14}\fancyfoot[C]{\thepage}\setlength{\headheight}{14pt}
\title{Burnside 引理与 Pólya 计数}\author{zhangjunshuo}\date{2026 年 10 月}
\begin{document}
\maketitle
\begin{quote}
本讲根据 2026 年 10 月 29 日的三页手写原稿、第二段课堂转写和课程教材整理。第一段转写不可用；手写稿第 1 页含 EGF 回顾（带标号二叉树、排列/轮换），转写中缺失，按手写稿与第 13 讲补全。整理版按“手写稿 + 转写 + 讲义/教材”取并集，不一致处给出勘误。

本讲主线：EGF 回顾 $\to$ 群作用与 Burnside 引理 $\to$ 循环指数与 Pólya 计数 $\to$ 无标号有根树的 Pólya 方程。

每个分节开头都先交代本节用到的记号；所有符号在首次使用处给出定义。
\end{quote}
\tableofcontents\clearpage

\section{从 EGF 到树与排列：本讲的前置例子}


这一讲的主线是\textbf{用群作用数“对称意义下不同的对象”}：Burnside 引理、循环指数与 Pólya 计数。开头的几个 EGF 例子是第 13 讲的延续，也正好提供了“带标号结构”的对照——后面无标号树要用普通生成函数与 Pólya 计数，不能混用。

\textbf{本节用到的记号：} $\hat A(x)=\sum a_nx^n/n!$ 为 EGF；$[x^n]$ 为系数提取；$b_n$ 为 $n$ 个结点的带标号树个数；$C_n$ 为循环群（这里不表示卡特兰数，按上下文区分）。

\subsection{EGF 的两条基本事实}


\textbf{乘积（二项卷积）.} 若 $\hat A,\hat B$ 分别是 $\{a_n\},\{b_n\}$ 的 EGF，则

\[
n!\,[x^n]\hat A(x)\hat B(x)=\sum_{i=0}^{n}\binom ni a_i b_{n-i}.
\]

组合意义：把 $n$ 个标签分成两部分，分别造 $\mathcal A$ 与 $\mathcal B$ 结构。

\textbf{集合构造.} 若“连通/不可分结构”的 EGF 是 $\hat G$，则由若干这样的结构\textbf{无序}拼成的对象的 EGF 是

\[
\exp\hat G=1+\hat G+\frac{\hat G^2}{2!}+\frac{\hat G^3}{3!}+\cdots
\]

（$\hat G^k$ 先给 $k$ 个有序部分，再除以 $k!$ 消去顺序。）

\subsection{例一：排列与轮换}


$n$ 元排列的 EGF 是

\[
\hat P(x)=\sum_{n\ge0}n!\frac{x^n}{n!}=\frac{1}{1-x}.
\]

一个 $k$ 元轮换有 $(k-1)!$ 个，故轮换的 EGF 是

\[
\hat C(x)=\sum_{k\ge1}(k-1)!\frac{x^k}{k!}
=\sum_{k\ge1}\frac{x^k}{k}=-\log(1-x).
\]

排列 = 轮换的无序集合，因此

\[
\exp\hat C(x)=\exp\bigl(-\log(1-x)\bigr)=\frac{1}{1-x}=\hat P(x),
\]

与“排列分解为轮换”一致。这条例子说明：\textbf{EGF 的 $\exp$ 对应集合构造}。

\subsection{例二：带标号的有根二叉树}


考虑每个结点至多有 $2$ 个\textbf{无序}孩子的有根树，结点带 $1,\dots,n$ 的标签；$b_n$ 为 $n$ 个结点的方案数，约定 $b_0=1$（空树）。

\textbf{递推/组合方程.} 去掉根后，剩下的孩子组成一个无序集合，大小可以是 $0,1,2$：

\begin{itemize}
\item $0$ 个孩子：$1$ 种；
\item $1$ 个孩子：一棵非空树，EGF 为 $T(x)=B(x)-1$；
\item $2$ 个孩子：两棵子树无序，EGF 为 $T(x)^2/2$。
\end{itemize}


再乘上根自己的标签 $x$，得

\[
T(x)=x\left(1+T(x)+\frac{T(x)^2}{2}\right),\qquad B(x)=1+T(x).
\]

消去 $T=B-1$，可写成

\[
B(x)=1+\frac{x}{2}\bigl(B(x)^2+1\bigr).
\]

\textbf{解方程.} 由 $xB^2-2B+(x+2)=0$，取 $x\to0$ 时 $B\to1$ 的分支：

\[
B(x)=\frac{1-\sqrt{1-2x-x^2}}{x}.
\]

展开得

\[
B(x)=1+x+x^2+\frac{3}{2}x^3+\frac{5}{2}x^4+\frac{9}{2}x^5+\cdots,
\]

所以

\[
b_0,b_1,b_2,b_3,b_4,b_5,\dots=1,1,2,9,60,540,\dots
\]

（$b_n=n![x^n]B$）。这些数可以用递推验证：$b_n$ 来自“根 + 无序的 $0/1/2$ 棵子树”。

\subsection{与下一节的衔接}


带标号树用 EGF，因为“标签”在乘法中要先分配；而\textbf{无标号}树（同构算一种）不能用 EGF 的 $\exp$，要用下一节的 Burnside/Pólya 或普通生成函数。这是本讲反复出现的分界。

\subsection{本节的提醒}


\begin{itemize}
\item \textbf{EGF 的 $x^n$ 系数是 $a_n/n!$}，读答案时要乘回 $n!$。
\item \textbf{$\exp$ 只对 EGF/带标号结构成立}；对普通生成函数直接写 $e^{A(x)}$ 一般没有组合意义。
\item 二叉树方程中“无序”二字很关键：有顺序孩子会得到不同的方程与计数。
\end{itemize}

\section{群作用、轨道与 Burnside 引理}


“对称意义下有多少种不同对象”这句话可以严格化为：群 $G$ 作用在集合 $S$ 上，数\textbf{轨道}的个数。Burnside 引理把轨道数化成“每个群元素的不动点个数”的平均，这是本讲所有计数公式的出发点。

\textbf{本节用到的记号：} $G\times S\to S$ 为左作用；$\operatorname{Orb}(x)=\{g\cdot x:g\in G\}$；$\operatorname{Stab}(x)=\{g\in G:g\cdot x=x\}$；$S/G$ 为轨道集合；$\operatorname{Fix}(g)=\{x\in S:g\cdot x=x\}$。

\subsection{群作用与轨道}


\textbf{定义.} 群 $G$ 在集合 $S$ 上的一个\textbf{左作用}是映射 $G\times S\to S$，$(g,x)\mapsto g\cdot x$，满足

\[
e\cdot x=x,\qquad (gh)\cdot x=g\cdot(h\cdot x).
\]

对每个 $g\in G$，$x\mapsto g\cdot x$ 是 $S$ 的一个置换；于是有同态 $G\to\operatorname{Sym}(S)$。

\textbf{轨道与稳定子.} $x$ 的轨道是它能被 $G$ 送到的一切元素；稳定子是固定 $x$ 的群元素。轨道是 $S$ 的划分，$S/G$ 表示全体轨道。

\textbf{命题（轨道--稳定子）.} 对任意 $x\in S$，

\[
|G|=|\operatorname{Orb}(x)|\cdot|\operatorname{Stab}(x)|,
\]

即 $|\operatorname{Orb}(x)|=[G:\operatorname{Stab}(x)]$。

\textbf{证明思路.} 构造双射 $G/\operatorname{Stab}(x)\to\operatorname{Orb}(x)$，$g\operatorname{Stab}(x)\mapsto g\cdot x$；良定义与单射来自稳定子的定义，满射显然。$\square$

\subsection{Burnside 引理}


\textbf{定理（Burnside）.} 有限群 $G$ 作用在有限集 $S$ 上，则轨道数为

\[
|S/G|=\frac{1}{|G|}\sum_{g\in G}|\operatorname{Fix}(g)|.
\]

\textbf{证明（双计数）.} 数集合

\[
X=\{(g,x)\in G\times S:g\cdot x=x\}.
\]

按 $g$ 数：$|X|=\sum_{g\in G}|\operatorname{Fix}(g)|$。按 $x$ 数：固定 $x$ 的 $g$ 有 $|\operatorname{Stab}(x)|$ 个，所以

\[
|X|=\sum_{x\in S}|\operatorname{Stab}(x)|
=\sum_{x\in S}\frac{|G|}{|\operatorname{Orb}(x)|}
=|G|\sum_{x\in S}\frac{1}{|\operatorname{Orb}(x)|}.
\]

同一轨道内的元素贡献相同，若轨道为 $O$，其中每个 $x$ 贡献 $1/|O|$，整条轨道共贡献 $1$。于是 $|X|=|G|\cdot|S/G|$，即得公式。$\square$

\textbf{使用要点.} 把 $S$ 取成“所有染色/标号方案”，$G$ 取成“允许的对称操作”；Burnside 数出的就是\textbf{对称意义下的方案数}。

\subsection{例一：四颗珠子的旋转染色}


用黑白两色给 $4$ 个位置染色，只允许旋转（群 $C_4$）。所有染色共 $2^4=16$ 种。

四个旋转的不动点数：

\begin{itemize}
\item 转 $0^\circ$（恒等）：$16$；
\item 转 $90^\circ$ 或 $270^\circ$：四颗必须同色，$2$ 种；
\item 转 $180^\circ$：两对相对位置分别同色，$2^2=4$ 种。
\end{itemize}


由 Burnside，

\[
|S/C_4|=\frac{1}{4}(16+2+4+2)=6.
\]

若进一步允许\textbf{翻转}（群 $D_4$），还要数 $4$ 个翻转的不动点：两条顶点轴各固定 $2^3=8$ 种，两条边轴各固定 $2^2=4$ 种，于是

\[
|S/D_4|=\frac{1}{8}(16+2+4+2+8+8+4+4)=6.
\]

（这个例子中两种群给出的轨道数恰好相同；一般并不相同。）

\subsection{例二：$GL_2(\mathbb F_q)$ 共轭作用}


$G=GL_2(\mathbb F_q)$ 通过 $g\cdot M=gMg^{-1}$ 作用在 $2\times2$ 矩阵集合 $M_2(\mathbb F_q)$ 上；轨道就是共轭类。由 Burnside，

\[
\#\text{共轭类}
=\frac{1}{|G|}\#\{(g,M):gM=Mg\}.
\]

按 $M$ 的中心化子大小分类可以算出结果为

\[
q^2+q.
\]

也可直接按特征多项式分类：标量矩阵 $q$ 类；非标量、有重特征值但不可对角化 $q$ 类；两个不同特征值（无序）$\binom q2$ 类；不可约二次特征多项式 $\frac{q^2-q}{2}$ 类，合计 $q^2+q$。这个例子说明 Burnside 也可以用来数“矩阵在相似意义下的类型”。

\subsection{本节的提醒}


\begin{itemize}
\item \textbf{$|\operatorname{Fix}(g)|$ 是对整个 $S$ 数的}，不是只数某个轨道。
\item \textbf{群作用要与“等价关系”一致}：旋转、翻转、置换标签分别对应不同的群；换群就换答案。
\item \textbf{轨道--稳定子}给出单个轨道大小；Burnside 给出轨道个数，两者不要混。
\item 恒等元的不动点数就是 $|S|$，可作为公式中的首项检查。
\end{itemize}

\section{循环指数与 Pólya 计数}


Burnside 只给出“颜色数固定时的方案数”。

若还想知道“用了几种颜色、每种颜色几个”，就要把每个群元素的\textbf{轮换结构}编码成多项式，再代入颜色权重——这就是循环指数与 Pólya 计数。

\textbf{本节用到的记号：} 置换 $g$ 的轮换型记为 $1^{c_1}2^{c_2}\cdots$，其中 $c_i(g)$ 是长度为 $i$ 的轮换个数。

$Z_G(t_1,t_2,\dots)$ 为循环指数；$C$ 为颜色数；$w_1,\dots,w_C$ 为颜色权重。

\subsection{置换的轮换型与循环指数}


\textbf{定义（循环指数）.} 有限群 $G$ 作用在有限集 $\Omega$ 上，对 $g\in G$ 记 $c_i(g)$ 为 $g$ 作为置换时长度为 $i$ 的轮换个数。定义

\[
Z_G(t_1,t_2,\dots)=\frac{1}{|G|}\sum_{g\in G}
t_1^{c_1(g)}t_2^{c_2(g)}\cdots .
\]

\textbf{例：循环群 $C_n$.} 旋转 $k$ 步的轮换个数是 $\gcd(k,n)$，每个轮换长度都是 $n/\gcd(k,n)$。令 $d=n/\gcd(k,n)$（即轮换长度），则 $d\mid n$，且满足 $\gcd(k,n)=n/d$ 的 $k$ 有 $\varphi(d)$ 个；每个这样的旋转贡献 $t_d^{n/d}$。因此

\[
Z_{C_n}(t_1,t_2,\dots)=\frac{1}{n}\sum_{d\mid n}\varphi(d)\,t_d^{\,n/d}.
\]

\textbf{例：二面体群 $D_n$（$2n$ 阶）.} 旋转部分同上；翻转部分分奇偶：

\begin{itemize}
\item $n$ 为奇数：$n$ 个翻转，每个有 $1$ 个不动点、$(n-1)/2$ 个对换，轮换型 $1^1 2^{(n-1)/2}$，贡献 $n\,t_1t_2^{(n-1)/2}$；
\item $n$ 为偶数：$n/2$ 个顶点轴翻转，轮换型 $1^2 2^{(n-2)/2}$，贡献 $\frac n2 t_1^2t_2^{(n-2)/2}$；$n/2$ 个边轴翻转，轮换型 $2^{n/2}$，贡献 $\frac n2 t_2^{n/2}$。
\end{itemize}


因此

\[
Z_{D_n}=\frac{1}{2n}\left(\sum_{d\mid n}\varphi(d)t_d^{n/d}
+\begin{cases}
n\,t_1t_2^{(n-1)/2}, & n\ \text{奇},\\[2pt]
\frac n2\left(t_1^2t_2^{(n-2)/2}+t_2^{n/2}\right), & n\ \text{偶}.
\end{cases}\right)
\]

\textbf{例（$n=4$）.}

\[
Z_{C_4}=\frac14\left(t_1^4+t_2^2+2t_4\right),\qquad
Z_{D_4}=\frac18\left(t_1^4+t_2^2+2t_4+2t_1^2t_2+2t_2^2\right).
\]

\subsection{Pólya 计数定理}


\textbf{定理（Pólya）.} 设颜色集为 $\{1,\dots,C\}$，颜色 $j$ 的权重为 $w_j$。用这些颜色给 $\Omega$ 染色，在 $G$ 作用下等价，则“按各颜色使用次数计权”的生成函数为

\[
Z_G\left(\sum_{j=1}^{C}w_j,\ \sum_{j=1}^{C}w_j^2,\ \sum_{j=1}^{C}w_j^3,\dots\right).
\]

若只数方案数（每种颜色权重取 $1$），令所有 $t_i=C$，得

\[
\#(\text{染色轨道})=Z_G(C,C,\dots).
\]

\textbf{为什么是 $t_i=\sum_j w_j^i$.} 在一个长度为 $i$ 的轮换中，置换会把位置循环移动，要保持染色不变，这 $i$ 个位置必须同色；选颜色 $j$ 时该轮换贡献 $w_j^i$。一个群元素 $g$ 的不动染色按轮换乘起来，正好是 $\prod_i(\sum_jw_j^i)^{c_i(g)}$；再用 Burnside 取平均即得。

\subsection{例：四颗珠子的双色染色}


取 $C_4$，颜色权重 $b,w$，代入 $t_i=b^i+w^i$：

\[
Z_{C_4}(b+w,b^2+w^2,b^4+w^4)
=\frac14\left((b+w)^4+(b^2+w^2)^2+2(b^4+w^4)\right).
\]

展开：

\[
b^4+b^3w+2b^2w^2+bw^3+w^4.
\]

各项系数就是“用 $0,1,2,3,4$ 颗黑珠”在旋转等价下的方案数（例如两黑两白有 $2$ 类：相邻与相对）。

若同时允许翻转（$D_4$），再代入 $Z_{D_4}$：

\[
\frac18\left[(b+w)^4+(b^2+w^2)^2+2(b^4+w^4)+2(b+w)^2(b^2+w^2)+2(b^2+w^2)^2\right].
\]

展开后也是 $b^4+b^3w+2b^2w^2+bw^3+w^4$——在 $n=4$ 这个例子中，旋转与翻转给出的计数多项式恰好相同；一般 $n$ 时两者不同。

\textbf{一般颜色数.} 令所有 $t_i=C$，由 $Z_{C_n}$ 得

\[
\#(\text{旋转等价})= \frac1n\sum_{d\mid n}\varphi(d)C^{n/d}.
\]

由 $Z_{D_n}$ 得

\[
\#(\text{翻转等价})=
\begin{cases}
\dfrac{1}{2n}\left(\sum_{d\mid n}\varphi(d)C^{n/d}+nC^{(n+1)/2}\right), & n\ \text{奇},\\[6pt]
\dfrac{1}{2n}\left(\sum_{d\mid n}\varphi(d)C^{n/d}+\dfrac n2\left(C^{n/2+1}+C^{n/2}\right)\right), & n\ \text{偶}.
\end{cases}
\]

\subsection{带“特殊颜色/不染色”的计数}


若要数“有些位置不染色，并关心不染色点的个数”，可把每个轮换的贡献改成 $C+x^i$（$C$ 种普通颜色，或用 $x$ 标记“这一个轮换不染色”），再代入循环指数。这样得到的多项式 $F(x)=\sum a_kx^k$ 中，$a_k$ 就是“恰有 $k$ 个点不染色”的轨道数。

\textbf{奇偶筛选.} 想只要不染色点为偶数个，取 $\frac{F(1)+F(-1)}2$；要 $3$ 的倍数等可用单位根/周期筛选。这是生成函数与 Burnside 结合的常见技巧。

\subsection{本节的提醒}


\begin{itemize}
\item \textbf{循环指数先看群，再看集合大小}：$Z_G$ 只依赖群在 $\Omega$ 上的置换表示。
\item \textbf{$D_n$ 的阶是 $2n$}：公式分母是 $2n$，不是 $n$；记号 $D_n$ 与 [DF] 的 $D_{2n}$ 相差下标。
\item \textbf{Pólya 代入的是 $t_i=\sum_j w_j^i$}，不是 $w_j$；指数 $i$ 来自轮换长度。
\item 颜色权重取 $1$ 时，$t_i=C$，才得到纯计数公式。
\end{itemize}

\section{应用：无标号有根树与 Pólya 方程}


这一节把 Burnside/Pólya 用到“无标号树”的计数：同构的树算一种。核心做法是把“根的孩子”看成一个\textbf{无序多重集}，用循环指数写生成函数方程。

\textbf{本节用到的记号：} $R(x)$ 为无标号有根树（每个结点至多 $3$ 个无序孩子；含空树）的普通生成函数；$R(x^k)$ 表示把 $x$ 换成 $x^k$；$x$ 记录结点数。

\subsection{为什么不能用 EGF}


带标号树的结点带 $1,\dots,n$ 的标签，$n$ 个结点的结构数用 EGF 很自然；无标号树把同构的树视为同一个，标签没了，\textbf{不能用 EGF 先分标签}。无标号结构要用普通生成函数，并用 Pólya 处理“孩子无序”的对称性。

\subsection{组合方程}


设 $\mathcal R$ 为“无标号有根树”类，每个结点至多有 $3$ 个无序孩子。记

\[
R(x)=\sum_{n\ge0}r_nx^n,
\]

其中 $r_0=1$ 代表\textbf{空树}（用于填充“空的孩子槽”），$r_n$ 为 $n$ 个结点的非空树个数。

\textbf{把每个结点看成“根 + 3 个孩子槽”.} 每个槽可以是空的（贡献 $1$）或一棵树（贡献 $R(x)$），三个槽\textbf{无序}。由 $S_3$ 的循环指数

\[
Z_{S_3}(t_1,t_2,t_3)=\frac16\left(t_1^3+3t_1t_2+2t_3\right),
\]

代入 $t_i=R(x^i)$（槽的轮换长度为 $i$ 时必须放同一棵树，故把 $x$ 换成 $x^i$），再乘根自己的 $x$，得到

\[
R(x)=1+\frac{x}{6}\left(R(x)^3+3R(x)R(x^2)+2R(x^3)\right).
\]

这里最前面的 $1$ 对应空树；右端的三个槽允许为空，所以“至多 $3$ 个孩子”被自动处理。

\textbf{另一种等价写法.} 若把 $R$ 定义为\textbf{非空}树的生成函数（常数项为 $0$），则方程写成

\[
R(x)=x\left(1+R(x)+\frac{R(x)^2+R(x^2)}{2}
+\frac{R(x)^3+3R(x)R(x^2)+2R(x^3)}{6}\right),
\]

两者只差“是否把空树算进 $R$”，系数相同。

\subsection{前几项}


把方程右端反复迭代（或按 $x$ 的幂逐项比较），得

\[
R(x)=1+x+x^2+2x^3+4x^4+8x^5+17x^6+39x^7+89x^8+\cdots
\]

即 $n=1,2,\dots$ 个结点的无标号有根树（每个结点至多 $3$ 个无序孩子）分别有

\[
1,1,2,4,8,17,39,89,\dots
\]

种。这个方程没有简单的闭式，通常用递推/迭代计算系数。

\subsection{与带标号情形的对照}


\begin{itemize}
\item \textbf{带标号、至多 2 个无序孩子}：EGF $T=x(1+T+T^2/2)$，计数 $1,2,9,60,540,\dots$（从 $n=1$ 起），乘法是二项卷积；
\item \textbf{无标号、至多 3 个无序孩子}：OGF 满足上面的 Pólya 方程，计数 $1,1,2,4,8,17,\dots$（从 $n=1$ 起），乘法是普通卷积 + 对称性修正。
\end{itemize}


两者都描述“孩子无序的树”，但一个带标签、一个不带标签，方程与数列完全不同；做题时先问“同构算不算同一种”。

\subsection{树的自同构与染色计数}


若给定一棵具体的树 $T$，想数它用 $C$ 种颜色的染色在 $\operatorname{Aut}(T)$ 作用下的轨道数，也可用 Burnside：

\[
\#(\text{染色轨道})=\frac{1}{|\operatorname{Aut}(T)|}
\sum_{\sigma\in\operatorname{Aut}(T)}C^{c(\sigma)},
\]

其中 $c(\sigma)$ 是 $\sigma$ 作为顶点置换的轮换个数。若还关心各颜色的使用次数，就把 $C^{c(\sigma)}$ 换成 $\left(\sum_j w_j\right)^{c(\sigma)}$ 的循环指数形式。这与 Pólya 公式是同一件事，只是群换成树的对称群。

\subsection{本节的提醒}


\begin{itemize}
\item \textbf{无标号与带标号不同}：不要对无标号树直接套 EGF 的 $\exp$。
\item \textbf{孩子无序时}，$2$ 棵与 $3$ 棵孩子的项要分别用 $S_2,S_3$ 的循环指数，不能只写 $R^2/2$、$R^3/6$（还要加 $R(x^2),R(x^3)$ 的修正项）。
\item \textbf{$R(x^k)$ 的系数不是 0/1 的简单重复}：它表示“同一棵子树重复 $k$ 次”的贡献，展开时按 $x$ 的幂合并。
\item 这类方程一般无闭式；考试若要求前几项，按系数逐项递推即可。
\end{itemize}


\section*{一页式总览}
\begin{itemize}
\item EGF：$\hat A(x)=\sum a_nx^n/n!$；乘法 $n![x^n]\hat A\hat B=\sum_i\binom ni a_ib_{n-i}$；$\exp$ 数无序集合。
\item 例：排列 $\hat P=1/(1-x)$、轮换 $\hat C=-\log(1-x)$、$\exp\hat C=\hat P$；带标号有根二叉树（孩子无序、至多 2 个）$T=x(1+T+T^2/2)$，$B=1+T=(1-\sqrt{1-2x-x^2})/x$，$b_n=1,1,2,9,60,540,\dots$。
\item 群作用：轨道 $\operatorname{Orb}(x)$、稳定子 $\operatorname{Stab}(x)$，$|G|=|\operatorname{Orb}(x)|\,|\operatorname{Stab}(x)|$。
\item Burnside：$|S/G|=\frac1{|G|}\sum_{g\in G}|\operatorname{Fix}(g)|$。
\item 四珠双色：旋转群 $C_4$ 下 $6$ 种；二面体群 $D_4$ 下也是 $6$ 种（分布 $b^4+b^3w+2b^2w^2+bw^3+w^4$）。
\item 循环指数：$Z_G(t_1,t_2,\dots)=\frac1{|G|}\sum_{g\in G}t_1^{c_1(g)}t_2^{c_2(g)}\cdots$。
\item $Z_{C_n}=\frac1n\sum_{d\mid n}\varphi(d)t_d^{n/d}$；$Z_{D_n}$ 的翻转项：$n$ 奇 $n\,t_1t_2^{(n-1)/2}$，$n$ 偶 $\frac n2(t_1^2t_2^{(n-2)/2}+t_2^{n/2})$。
\item Pólya：把 $t_i$ 换成 $\sum_j w_j^i$；颜色数 $C$ 时 $t_i=C$，得染色轨道数。
\item 项链数：$C_n$ 群下 $\frac1n\sum_{d\mid n}\varphi(d)C^{n/d}$；$D_n$ 群下按奇偶加翻转项再除 $2n$。
\item 无标号树：$R(x)=1+\frac{x}{6}(R(x)^3+3R(x)R(x^2)+2R(x^3))$；$R=1+x+x^2+2x^3+4x^4+8x^5+17x^6+\cdots$。
\end{itemize}

\section*{参考资料}
\begin{itemize}
\item \textbf{[DS]} 邓小铁、王畅。《离散数学与结构》第 14 章：§14.1 群作用和计数原理、§14.2 Burnside 引理、§14.3 Pólya 定理。
\item \textbf{[讲义]} 课程讲义 \texttt{25-26-lecture\_notes.pdf} §9.10 Burnside's Lemma、§9.11 Polya Counting、§9.12 共轭类计数。
\item 原始材料：本地归档（手写原稿、第二段转写）。
\end{itemize}
\end{document}
